One-shot federated distillation combines local teachers into a global student, but class specialization makes their contributions uneven. We study whether selecting teachers by measured class competence adds value beyond knowing which classes they encountered during training, and whether restricting selected teachers' outputs improves the distilled student. Our protocol uses a labeled public proxy and separate private training, validation and expertise splits. Across three datasets, six client allocations and three seeds, we compare selection rules, logit and probability pooling, support restriction, and proxy-only supervision. The agreed study contains 471 student runs. Under disjoint two-class and single-class allocations, presence and expertise masks coincide in the observed conditions: large improvements over uniform pooling do not establish an additional benefit from measuring competence there. Under CIFAR-10 Dirichlet allocation with \(\alpha\)=0.5, measured competence adds 5.25 percentage points of mean accuracy over presence, while other regimes expose accuracy--likelihood trade-offs. Restriction lowers target negative log-likelihood in all six CIFAR-10 regimes but increases student negative log-likelihood in all six and reduces student accuracy in five. Finally, a proxy-size study shows that improvements over uniform distillation need not imply improvements over direct supervision. These results distinguish support-aware selection from competence accreditation and identify limits to inferring student utility from target quality.
